\chapter{プロンプトエンジニアリング}
\label{chap2}

\begin{chapterbox}
\textbf{【この章で学ぶこと】}
\begin{itemize}
\item 効果的なプロンプトの設計原則（ROCFモデル）を理解する
\item Zero-shot、Few-shot、Chain-of-Thoughtなどの技法を使い分けられる
\item ハルシネーションを検出し対処する方法を身につける
\end{itemize}
\end{chapterbox}
\vspace{5mm}

%=========================================
\section{なぜプロンプトが重要か}
\label{sec:2.1}
%=========================================

\subsection{プロンプト次第で結果は大きく変わる}

生成AIから良い出力を得るためには、良い入力（プロンプト\index{ぷろんぷと@プロンプト}）が欠かせない。同じ目的であっても、プロンプトの書き方によって、得られる回答の質は大きく異なる。

以下の例を見てほしい。ある修士学生が、自分の研究テーマに関する先行研究を調べたい場面を想定する。

\begin{promptbox}
\textbf{【悪いプロンプトの例】}

深層学習の論文を教えて
\end{promptbox}

この質問に対して、AIは「深層学習」という広大な分野の中から、何を回答すべきか判断できない。結果として、一般的すぎる回答や、質問者の意図とずれた回答が返ってくることが多い。

\begin{promptbox}
\textbf{【良いプロンプトの例】}

私は修士1年で、「Few-shot学習を用いた製造業における\\
外観検査の自動化」をテーマに研究を始めます。

以下の条件で、関連する先行研究を5本紹介してください。

条件:
\begin{itemize}
\item 2021年以降に発表された査読付き論文
\item Few-shot learningまたはMeta-learningの手法を用いている
\item 製造業の品質管理に関連する応用研究
\item 各論文について、タイトル、著者、掲載誌、主要な貢献を\\
100字程度で説明してください
\end{itemize}

注意: 実在する論文のみを挙げてください。確信が持てない場合は「確認が必要」と明記してください。
\end{promptbox}

このプロンプトでは、背景（修士1年で研究を始める段階）、具体的な研究テーマ、求める情報の条件、出力形式が明確に指定されている。AIはこれらの情報を手がかりに、より適切な回答を生成できる。

\begin{screen}
\textbf{キーポイント}: プロンプトは「AIへの仕様書」である。ソフトウェア開発において曖昧な仕様書から良い製品が生まれないのと同様に、曖昧なプロンプトから良い回答は生まれない。
\end{screen}

\subsection{プロンプトエンジニアリングとは}

\textbf{プロンプトエンジニアリング}\index{ぷろんぷとえんじにありんぐ@プロンプトエンジニアリング}（Prompt Engineering）とは、AIから望ましい出力を得るために、入力（プロンプト）を設計・最適化する技術のことである\cite{liu2023pretrain}。

プロンプトエンジニアリングは、単なる「質問の仕方」にとどまらない。それは以下の要素を含む、体系的な設計プロセスである。

\begin{enumerate}
\item \textbf{目的の明確化}: AIに何をさせたいのかを明確にする
\item \textbf{文脈の設計}: AIが適切な回答を生成するために必要な背景情報を選定する
\item \textbf{制約の指定}: 出力の形式、長さ、スタイルなどの条件を設定する
\item \textbf{反復的改善}: 結果を評価し、プロンプトを修正するサイクルを回す
\end{enumerate}

研究者にとってプロンプトエンジニアリングは、今後ますます重要なスキルとなる。それは、AIの能力を引き出すためにも、AIの限界を見極めるためにも重要である。

%=========================================
\section{ROCFモデル：良いプロンプトの4要素}
\label{sec:2.2}
%=========================================

本書では、効果的なプロンプトを設計するためのフレームワークとして、\textbf{ROCFモデル}\index{ROCFもでる@ROCFモデル}を提案する。ROCFは以下の4要素の頭文字である。

\begin{itembox}[l]{ROCFモデルの4要素}
\begin{itemize}
\item \textbf{R}ole（役割）--- AIに専門家の役割を与える
\item \textbf{O}bjective（目的）--- 何をしてほしいか明確に
\item \textbf{C}ontext（文脈）--- 背景情報を提供する
\item \textbf{F}ormat（形式）--- 出力形式を指定する
\end{itemize}
\end{itembox}

\subsection{Role（役割）：AIに専門家の役割を与える}

AIに特定の専門家の役割を付与することで、その分野に適した語彙、視点、判断基準で回答を生成させることができる。

\begin{promptbox}
\textbf{【Role要素の例】}

例1: ``あなたは計算機科学分野の経験豊富な査読者です。''

例2: ``あなたは統計学の教授で、大学院生の研究指導を行っています。''

例3: ``あなたはテクニカルライターで、学術論文の英語校正を専門としています。''
\end{promptbox}

Roleの効果は、単に回答のトーンを変えるだけではない。たとえば、「査読者として」と指定すると、AIは論文の弱点や改善点を指摘する方向で回答を生成する傾向がある。「指導教員として」と指定すると、学生の理解を助けるような丁寧な説明を生成しやすくなる。

\begin{screen}
\textbf{ヒント}: Roleは具体的であるほど効果的である。「専門家として」よりも「機械学習分野で20年の研究経験を持つ教授として、修士学生の論文を指導する立場で」と指定した方が、より適切な回答を得やすい。
\end{screen}

\subsection{Objective（目的）：何をしてほしいか明確に}

Objectiveは、AIに実行してほしいタスクを明確に指定する要素である。曖昧な依頼は曖昧な結果を生む。

\begin{promptbox}
\textbf{【Objective要素の比較】}

曖昧な目的: ``この論文について教えてください''

明確な目的: ``この論文の研究手法と実験結果を、\\
以下の3つの観点から分析してください:\\
(1) 新規性、(2) 再現性、(3) 既存手法との比較''
\end{promptbox}

Objectiveを書く際のポイント:
\begin{enumerate}
\item \textbf{動詞を明確にする}: 「教えてください」ではなく「要約してください」「比較してください」「評価してください」など
\item \textbf{範囲を限定する}: 「すべて」ではなく「上位3つ」「主要な点」「特にXに関して」など
\item \textbf{成功基準を示す}: 何をもって良い回答とするかを伝える
\end{enumerate}

\subsection{Context（文脈）：背景情報を提供する}

Contextは、AIが適切な回答を生成するために必要な背景情報である。研究活動では特に重要な要素となる。

\begin{promptbox}
\textbf{【Context要素の例】}

``以下の背景を踏まえて回答してください:
\begin{itemize}
\item 私は情報工学専攻の修士1年で、深層学習を用いた自然言語処理が専門です
\item 対象とする言語は日本語です
\item 使用可能な計算資源はGPU 1台（NVIDIA A100）です
\item 投稿先はACL 2027を予定しています
\item 研究期間は1年半です''
\end{itemize}
\end{promptbox}

適切なContextを提供することで、AIはより現実的で実行可能な提案を行えるようになる。たとえば、「GPU 1台」という制約を伝えれば、計算コストが膨大なモデルの提案を避け、効率的なアプローチを推薦してくれる。

\subsection{Format（形式）：出力形式を指定する}

Formatは、AIの出力がどのような形式であるべきかを指定する要素である。

\begin{promptbox}
\textbf{【Format要素の例】}

例1: ``箇条書きで5項目にまとめてください''

例2: ``表形式で比較してください（列: 手法名、精度、計算コスト、利点、欠点）''

例3: ``\LaTeX 形式で数式を含めて記述してください''

例4: ``以下のJSON形式で出力してください''
\end{promptbox}

\subsection{ROCFモデルの実践——4要素を組み合わせる}

4つの要素を組み合わせた実践的なプロンプトの例を示す。

\begin{promptbox}
\textbf{【ROCFモデルの完全な適用例】}

\textbf{[Role]} あなたは自然言語処理分野の研究に精通した\\
大学教授です。修士学生の研究指導を日常的に行っています。

\textbf{[Objective]} 以下の研究テーマについて、修士論文の\\
構成案（章立て）を提案してください。各章の概要と、\\
そこで扱うべき内容を200字程度で説明してください。

\textbf{[Context]}
\begin{itemize}
\item 研究テーマ: 「事前学習済み言語モデルを用いた日本語法律文書の自動要約」
\item 修士2年の学生で、修士論文の提出は6か月後です
\item すでにベースラインの実験は完了し、提案手法の実装中です
\item 関連研究として、BERTを用いた法律文書分類の研究はすでにサーベイ済みです
\end{itemize}

\textbf{[Format]} 以下の形式で出力してください:

第X章 [章タイトル] / 概要: [200字程度の説明] / 含めるべき内容: [項目リスト]
\end{promptbox}

\begin{screen}
\textbf{キーポイント}: ROCFモデルの4要素すべてを毎回使う必要はない。タスクの性質に応じて、必要な要素を選択すればよい。ただし、ObjectiveとContextは常に含めることを推奨する。
\end{screen}

%=========================================
\section{基本的なプロンプト技法}
\label{sec:2.3}
%=========================================

\subsection{Zero-shot Prompting（例なしで指示）}

\textbf{Zero-shot prompting}\index{Zero-shot prompting}は、具体例を一切提示せず、指示文のみでタスクを実行させる技法である。最も基本的なプロンプト技法であり、明確なタスクには十分に機能する。

\begin{promptbox}
\textbf{【Zero-shot promptingの例】}

以下の論文タイトルを、学術的な日本語に翻訳してください。

``Retrieval-Augmented Generation for Knowledge-Intensive NLP Tasks''
\end{promptbox}

\textbf{適している場面:}
\begin{itemize}
\item タスクが明確で、AIが十分な学習データを持っている場合
\item 翻訳、要約、校正などの定型的なタスク
\item 簡単な質問応答
\end{itemize}

\textbf{注意点:}
\begin{itemize}
\item タスクの定義が曖昧な場合、AIが意図と異なる解釈をする可能性がある
\item 出力形式を厳密に制御したい場合は、追加の指定が必要
\end{itemize}

\subsection{Few-shot Prompting（例を添えて指示）}

\textbf{Few-shot prompting}\index{Few-shot prompting}は、望ましい入出力の例を数個提示した上で、タスクを実行させる技法である\cite{brown2020language}。AIに「こういう形式で、こういう品質の回答を期待している」というお手本を示すことで、出力の品質と一貫性が大幅に向上する。

\begin{promptbox}
\textbf{【Few-shot promptingの例：論文の分類】}

以下の論文タイトルを、研究分野に分類してください。

例1:

タイトル: ``BERT: Pre-training of Deep Bidirectional Transformers for Language Understanding''

分野: 自然言語処理 $>$ 事前学習モデル

例2:

タイトル: ``ResNet: Deep Residual Learning for Image Recognition''

分野: コンピュータビジョン $>$ 画像認識

例3:

タイトル: ``Playing Atari with Deep Reinforcement Learning''

分野: 強化学習 $>$ ゲームAI

では、以下の論文を分類してください:

タイトル: ``Generative Adversarial Nets''

分野:
\end{promptbox}

\textbf{適している場面:}
\begin{itemize}
\item 出力の形式や品質を厳密に制御したい場合
\item AIがあまり得意でないタスク（特殊な分類基準など）
\item 独自の評価基準や記述スタイルを適用したい場合
\end{itemize}

\textbf{Few-shotの例数の目安:}
\begin{itemize}
\item 2〜5個の例があれば、多くのタスクで十分
\item 例が多すぎると、コンテキスト長を浪費する
\item 例は多様性を持たせる（似た例ばかりにしない）
\end{itemize}

\subsection{Chain-of-Thought（ステップバイステップの思考を促す）}

\textbf{Chain-of-Thought（CoT）prompting}\index{Chain-of-Thought}は、AIに思考過程を段階的に出力させる技法である\cite{Wei2022ChainOfThought}。複雑な推論が必要なタスクで、回答の正確性が著しく改善される。

\begin{promptbox}
\textbf{【Chain-of-Thought promptingの例：研究手法の選択】}

以下の研究課題に対して、最適な実験手法を選択してください。\\
ステップバイステップで考え、各ステップの根拠を示してください。

研究課題: 「SNS投稿のテキストから、投稿者の感情状態を\\
4カテゴリ（喜び、怒り、悲しみ、恐怖）に分類するモデルを\\
開発したい。ただし、ラベル付きデータは各カテゴリ100件\\
しかない。」

ステップ1: まず、この課題の本質的な要件を整理してください

ステップ2: 考えられる手法を3つ以上列挙してください

ステップ3: 各手法のメリット・デメリットを分析してください

ステップ4: データ量の制約を考慮して、最適な手法を推薦してください

ステップ5: 推薦した手法の実装計画を簡潔に示してください
\end{promptbox}

Chain-of-Thoughtの効果が特に高い場面:
\begin{itemize}
\item 数学的な推論が必要な場合
\item 複数の要因を考慮した判断が必要な場合
\item 複雑な比較・分析が必要な場合
\end{itemize}

\begin{screen}
\textbf{キーポイント}: 「ステップバイステップで考えてください」という一文を追加するだけでも、回答の質が向上することが多い。これは、AIが中間的な推論過程を生成することで、最終的な結論の正確性が高まるためである。
\end{screen}

\subsection{各技法の比較と使い分け}

\begin{table}[htbp]
\centering
\caption{プロンプト技法の比較}
\label{tab:prompt_techniques}
\small
\begin{tabular}{llcl}
\toprule
技法 & 適用場面 & 例の必要数 & 効果 \\
\midrule
Zero-shot & 明確な定型タスク & 0 & 速い、簡潔 \\
Few-shot & 形式の制御が必要 & 2〜5 & 出力品質の安定化 \\
Chain-of-Thought & 複雑な推論 & 0〜3 & 推論精度の向上 \\
\bottomrule
\end{tabular}
\end{table}

実際の研究活動では、これらの技法を組み合わせて使うことが多い。たとえば、Few-shotの例にChain-of-Thoughtの思考過程を含めることで、「こういう思考プロセスで、こういう形式の回答を期待している」ということをAIに伝えることができる（Few-shot CoT）。

%=========================================
\section{高度なプロンプト技法}
\label{sec:2.4}
%=========================================

基本的な技法を習得したら、次はより高度な技法に挑戦しよう。これらの技法は、研究活動の中でも特に複雑なタスクで効果的である。

\subsection{ペルソナ設定}

ペルソナ設定\index{ぺるそなせってい@ペルソナ設定}は、ROCFモデルのRole要素をさらに発展させた技法である。AIに特定の視点から回答させることで、多角的な分析や批判的なフィードバックを得ることができる。

\begin{promptbox}
\textbf{【ペルソナ設定の実践例：論文の多角的レビュー】}

\textbf{--- ペルソナ1: 厳格な査読者 ---}

あなたは、トップカンファレンス（NeurIPS, ICML）の査読を\\
10年以上務めてきた厳格な研究者です。以下の論文概要に対して、\\
主要な弱点と改善すべき点を5つ指摘してください。\\
特に、技術的な新規性と実験の妥当性について厳しく評価してください。

{[}論文概要をここに貼る]
\end{promptbox}

\begin{promptbox}
\textbf{--- ペルソナ2: 好意的な指導教員 ---}

あなたは、修士学生の研究指導に情熱を持つ教授です。\\
以下の論文概要に対して、良い点を3つ、改善の余地がある点を\\
3つ指摘し、各改善点に対して具体的なアドバイスを提供して\\
ください。学生のモチベーションを維持しながら、\\
建設的なフィードバックを心がけてください。

{[}論文概要をここに貼る]
\end{promptbox}

\begin{promptbox}
\textbf{--- ペルソナ3: 隣接分野の研究者 ---}

あなたは、情報検索（IR）分野の研究者です。\\
自然言語処理（NLP）の専門家ではありませんが、テキスト処理に\\
関連する知識を持っています。以下の論文概要を読んで、\\
IR分野の研究者として疑問に思う点や、\\
IR分野との接点を指摘してください。

{[}論文概要をここに貼る]
\end{promptbox}

同じ論文概要に対して、異なるペルソナからのフィードバックを得ることで、自分では気づかなかった観点や問題点を発見できることがある。

\subsection{段階的精緻化（Iterative Refinement）}

段階的精緻化\index{だんかいてきせいちか@段階的精緻化}は、AIとの対話を複数ターンに分け、出力を少しずつ改善していく技法である。一度で完璧な出力を求めるのではなく、反復的に品質を高めていく。

\begin{promptbox}
\textbf{【段階的精緻化の例：研究要旨の作成】}

\textbf{--- ターン1: ドラフト生成 ---}

以下の情報をもとに、国際会議に投稿する論文のAbstract（250語以内、英語）のドラフトを作成してください。

研究テーマ: [テーマ] / 手法: [手法の概要] / 結果: [主要な結果] / 貢献: [主要な貢献]

\textbf{--- ターン2: 構造の改善 ---}

ありがとうございます。以下の点を修正してください:
(1) 最初の文で、研究の動機をより明確にしてください
(2) 手法の説明をもう少し具体的にしてください
(3) 結果の数値を追加してください: 精度92.3\%, ベースライン比+4.7ポイント

\textbf{--- ターン3: 言語の洗練 ---}

構造は良くなりました。次に、以下の観点で言語面を改善してください:
(1) 受動態を能動態に変えられる箇所を修正
(2) 冗長な表現を簡潔にする
(3) 学術的な語彙として不適切な表現があれば修正
\end{promptbox}

段階的精緻化のメリット:
\begin{itemize}
\item 各ステップで焦点を絞れるため、改善の方向が明確になる
\item AIの出力を確認しながら修正指示を出せるため、最終的な品質が高くなる
\item 人間の意図をAIに徐々に正確に伝えることができる
\end{itemize}

\subsection{メタプロンプティング}

\textbf{メタプロンプティング}\index{めたぷろんぷてぃんぐ@メタプロンプティング}は、AIに「プロンプト自体を考えさせる」技法である\cite{zhou2023large}。自分のタスクに最適なプロンプトがわからない場合に有効である。

\begin{promptbox}
\textbf{【メタプロンプティングの例】}

私は以下のタスクを実行したいのですが、最適なプロンプトの書き方がわかりません。

タスク: 修士論文の「関連研究」セクションを書くために、10本の論文を体系的に整理し、比較分析したい。

このタスクに対して、最も効果的なプロンプトを設計してください。以下の点を含めてください:
\begin{enumerate}
\item プロンプトの全文
\item なぜその書き方が効果的なのかの説明
\item プロンプトをさらに改善するためのヒント
\end{enumerate}
\end{promptbox}

メタプロンプティングは、プロンプトエンジニアリングの学習そのものにも有効である。AIが設計するプロンプトを分析することで、効果的なプロンプトの特徴を学ぶことができる。

\subsection{制約条件の活用}

AIの出力を精密に制御するためには、適切な制約条件\index{せいやくじょうけん@制約条件}を設定することが有効である。

\begin{promptbox}
\textbf{【制約条件の活用例】}

以下の制約条件に従って、研究計画書のドラフトを作成してください。

制約条件:
\begin{enumerate}
\item 文字数: 2000字以内
\item 構成: 背景、目的、手法、期待される成果の4セクション
\item 文体: 「である」調、学術的な表現
\item 含めるべきキーワード: 深層学習、転移学習、ドメイン適応
\item 避けるべき表現: 「画期的」「世界初」などの誇張表現
\item 参考にすべきスタイル: 科研費申請書の研究計画調書
\item 対象読者: 情報工学分野の審査委員（AI専門とは限らない）
\end{enumerate}
\end{promptbox}

効果的な制約条件のカテゴリを表\ref{tab:constraint_categories}に示す。

\begin{table}[htbp]
\centering
\caption{制約条件のカテゴリと例}
\label{tab:constraint_categories}
\begin{tabular}{ll}
\toprule
カテゴリ & 例 \\
\midrule
長さの制約 & 「300字以内」「5項目で」 \\
形式の制約 & 「表形式で」「\LaTeX で」「箇条書きで」 \\
内容の制約 & 「技術的な詳細は省く」「数式を含める」 \\
否定の制約 & 「専門用語を使わない」「主観的な評価は含めない」 \\
品質の制約 & 「学術論文レベルの正確さで」「初学者にもわかるように」 \\
スタイルの制約 & 「である調」「フォーマルな英語」 \\
\bottomrule
\end{tabular}
\end{table}

\begin{screen}
\textbf{ヒント}: 制約条件は多すぎると逆効果になることがある。AIが全ての制約を同時に満たそうとして、かえって不自然な出力になる場合がある。重要度の高い制約を3〜5個に絞ることを推奨する。
\end{screen}

%=========================================
\section{実践的なプロンプト設計技法}
\label{sec:practical_prompt}
%=========================================

前節までに紹介した技法に加え、AIサービス提供企業が公開しているベストプラクティスから、研究活動に役立つ実践的な技法を紹介する。本節では、特にAnthropicのプロンプトエンジニアリングガイドに基づく手法を取り上げる。

\subsection{XMLタグによるプロンプトの構造化}

複雑なプロンプトを作成する際、\textbf{XMLタグ}\index{XMLたぐ@XMLタグ}を用いてプロンプトの各部分を明示的に区切ることで、AIが指示・文脈・入力データを曖昧さなく解釈できるようになる。たとえば、指示部分を\texttt{<instructions>}タグで、分析対象の文書を\texttt{<document>}タグで囲むといった方法である。

この技法は、指示文、背景情報、具体的なデータ、出力形式の指定など、複数の異なる種類の情報を1つのプロンプトに含める場合に特に有効である。タグによって各要素の境界が明確になるため、AIが指示の一部をデータと誤認したり、データの一部を指示として解釈するといった問題を防ぐことができる。

研究活動では、論文をAIに分析させる場面でこの技法が威力を発揮する。論文テキストを\texttt{<document>}タグで囲み、分析の依頼を別のタグで明示することで、AIは「何が分析対象で、何が指示なのか」を正確に把握できる。

\begin{promptbox}
\textbf{【XMLタグを用いた構造化プロンプトの例】}

\texttt{<instructions>}

以下の論文のアブストラクトを読み、背景・手法・結果・限界の4点で要約してください。

\texttt{</instructions>}

\texttt{<document>}

{[}論文のアブストラクトをここに貼り付け]

\texttt{</document>}

\texttt{<output\_format>}

箇条書きで、各項目2--3文程度にまとめてください。

\texttt{</output\_format>}
\end{promptbox}

\subsection{長文コンテキストの効果的な配置}
\index{こんてきすとはいち@コンテキスト配置}

長い文書（論文全文やデータセット）をAIに与えて分析や質問応答を行う場合、\textbf{文書をプロンプトの先頭に配置し、質問や指示を末尾に配置する}ことで、回答品質が向上することがAnthropicの検証により報告されている。改善幅は最大30\%に達する場合もあるとされる。

この配置が効果的な理由は、LLMがプロンプト末尾の情報に対してより強い注意を向ける傾向があるためである。長い文書の後に質問を配置することで、AIは「いま何を求められているのか」を明確に認識した状態で文書内容を参照できる。

研究活動では、論文の全文をAIに読み込ませて質問する場面が多い。そのような場合、「論文テキスト→質問」の順にプロンプトを構成すると、回答の的確さが向上する。

\begin{screen}
\textbf{ヒント}: 長い論文やデータをAIに分析させるとき、文書を先に、質問を後に配置すると回答品質が向上する。「資料→指示」の順を心がけよう。
\end{screen}

\subsection{引用ベースの回答で正確性を高める}
\index{いんようべーすかいとう@引用ベース回答}

AIに文書の分析を依頼する際、まず\textbf{ソース文書から関連部分を引用させ、その後に分析を述べさせる}という手順を踏ませることで、ハルシネーションを大幅に抑制できる。この手法は\textbf{グラウンディング}\index{ぐらうんでぃんぐ@グラウンディング}（grounding）と呼ばれ、AIの回答を原文に「根拠づける」効果がある。

引用を先に行わせることで、AIは回答の根拠を原文に求めるようになり、「もっともらしいが原文にない情報」を生成するリスクが低減する。これは文献レビューや論文査読のように、原文の正確な参照が求められる研究場面で特に有用である。

\begin{promptbox}
\textbf{【引用ベースの回答を促すプロンプトの例】}

以下の論文から、研究手法に関する記述を引用してください。

引用部分を\texttt{<quotes>}タグで囲み、その後に\texttt{<analysis>}タグ内で分析を行ってください。

{[}論文テキスト]
\end{promptbox}

\subsection{指示の明確化：「やるな」より「代わりにこうして」}

AIに出力を制御する指示を与える際、\textbf{否定形の指示（「〜しないでください」）よりも、肯定形の指示（「代わりに〜してください」）の方が効果的}である。否定形の指示は「何をすべきでないか」を伝えるだけで、AIに明確な行動の方向を示さない。一方、肯定形の指示はAIが取るべき行動を具体的に示すため、出力がより安定する。

\vspace{3mm}
\textbf{悪い例と良い例の比較:}
\vspace{2mm}

\begin{itemize}
\item \textbf{悪い例}: 「箇条書きにしないでください」

$\rightarrow$ AIは箇条書き以外の何をすべきか判断に迷う

\item \textbf{良い例}: 「段落形式で、各段落に1つのポイントを含めて書いてください」

$\rightarrow$ AIは具体的な出力形式を把握できる
\end{itemize}

\vspace{2mm}

もう一つの例を挙げる。

\begin{itemize}
\item \textbf{悪い例}: 「マークダウンを使わないでください」
\item \textbf{良い例}: 「滑らかな段落構成の文章で書いてください」
\end{itemize}

\subsection{文脈と理由を添える}

AIへの指示では、\textbf{何をしてほしいかだけでなく、なぜそうしてほしいかを説明する}ことで、AIがその意図を汎化的に理解し、より適切に対応できるようになる。理由が明確であれば、指示で明示的にカバーしていない状況でも、AIは意図に沿った判断を下しやすくなる。

\vspace{3mm}
\textbf{悪い例と良い例の比較:}
\vspace{2mm}

\begin{itemize}
\item \textbf{悪い例}: 「省略記号（\ldots）を使わないでください」

$\rightarrow$ AIは指示には従えるが、なぜかわからないため類似の状況で判断を誤る可能性がある

\item \textbf{良い例}: 「この文章はテキスト読み上げソフトで使用するため、省略記号（\ldots）は使わないでください。読み上げソフトが省略記号を正しく発音できないためです。」

$\rightarrow$ AIは理由を理解し、省略記号以外にも読み上げに問題がありそうな表現を自主的に避けられる
\end{itemize}

\vspace{2mm}

研究活動では、たとえば「数式は\LaTeX 形式で書いてください。この文章をOverleafにそのまま貼り付けて使うためです。」のように理由を添えることで、数式以外の部分でも\LaTeX に適した記法をAIが選択するようになる。

\begin{screen}
本節の内容はAnthropicのプロンプトエンジニアリングガイド（Anthropic公式サイト「Prompting best practices」）を参考にしている。各AIツールの最新のベストプラクティスは、公式ドキュメントを確認されたい。
\end{screen}

%=========================================
\section{ハルシネーションの理解と対策}
\label{sec:2.5}
%=========================================

\subsection{ハルシネーションとは何か}

\textbf{ハルシネーション}\index{はるしねーしょん@ハルシネーション}（Hallucination）とは、AIがもっともらしいが事実に反する情報を生成する現象である。第\ref{chap1}章でも触れたが、プロンプトエンジニアリングの観点から、ここでより詳しく解説する。

ハルシネーションが生じる技術的な背景を簡潔に説明する。LLMは「次のトークン（単語）を予測する」モデルである。この予測は、学習データから得られた統計的パターンに基づいている。つまり、LLMは「事実を記憶している」のではなく、「もっともらしい文を生成している」のである。

この「もっともらしさ」と「正確さ」のギャップがハルシネーションの本質である。

\subsection{研究活動で特に危険なハルシネーション}

研究活動では、以下のタイプのハルシネーションが特に危険である。

\paragraph{タイプ1: 架空の論文引用}
最も頻繁に報告されるハルシネーションの一つである。AIは、実在しない論文のタイトル、著者名、雑誌名、出版年、さらにはDOI番号までもっともらしく生成することがある。

\begin{promptbox}
\textbf{【架空の論文引用の例】}

質問: ``ドメイン適応における最新の手法について、主要な論文を3本教えてください''

AIの回答（ハルシネーションを含む）:

1. Zhang, W., \& Liu, H. (2024). ``Progressive Domain\\
Adaptation with Contrastive Learning for Visual Recognition.''\\
IEEE TPAMI, 46(3), 1234--1248.

$\rightarrow$ この論文は実在しない可能性がある。著者名、タイトル、\\
巻号ページがもっともらしく生成されているが、\\
DOIで確認すると存在しない。
\end{promptbox}

\paragraph{タイプ2: 誤った数値データ}\mbox{}

\begin{promptbox}
\textbf{【誤った数値データの例】}

質問: ``ImageNetにおける最新の画像認識精度のSOTAは？''

AIの回答: ``2025年時点で、ImageNetにおける\\
Top-1精度のSOTAは94.7\%であり、\\
これはXYZモデルによって達成されました。''

$\rightarrow$ 数値やモデル名が不正確である可能性がある。\\
特にSOTAは頻繁に更新されるため、AIの学習データの\\
カットオフ以降の情報は誤っている可能性が高い。
\end{promptbox}

\paragraph{タイプ3: もっともらしい虚偽の説明}\mbox{}

\begin{promptbox}
\textbf{【もっともらしい虚偽の説明の例】}

質問: ``Batch Normalizationが有効な理由を教えてください''

AIの回答: ``Batch Normalizationは、内部共変量シフト\\
（Internal Covariate Shift）を軽減することで\\
学習を安定化させます。''

$\rightarrow$ 長年信じられてきた説明だが、\\
Santurkar et al.\ (2018)\cite{santurkar2018batch}の\\
研究により、\\
この説明が主要因ではないことが示されている。\\
AIは「広く引用されている説明」を生成しがちであり、\\
最新の理解を反映していない可能性がある。
\end{promptbox}

\subsection{5つの対策法}

ハルシネーションに対処するための5つの具体的な対策法を紹介する。

\begin{itembox}[l]{ハルシネーション対策の5つの方法}
\textbf{対策1: 出典を要求する}

プロンプトの中で、明示的に出典の提示を要求する。また、出典の信頼性を評価する基準も指定する。

\vspace{2mm}
\textbf{対策2: 複数ツールでクロスチェック}

重要な情報は、複数のAIツールや検索エンジンで確認する。

\vspace{2mm}
\textbf{対策3: 原典にあたって確認}

AIが提示した情報の元となる原典（論文、公式ドキュメント、データソース）を直接確認する。これは最も確実な対策であり、研究者の基本動作でもある。

\vspace{2mm}
\textbf{対策4: 「わからない」と言わせる}

AIに「わからない場合はわからないと回答してよい」と明示的に許可することで、ハルシネーションを抑制できることがある。

\vspace{2mm}
\textbf{対策5: 段階的に検証する}

AIの回答を一度に全て信頼するのではなく、小さな単位に分割して段階的に検証する。
\end{itembox}

\begin{promptbox}
\textbf{【出典を要求するプロンプト例】}

以下の質問に回答する際、各主張に対して出典を明記してください。出典は以下の形式で記載してください:
\begin{itemize}
\item 著者名、論文タイトル、掲載誌/会議名、年
\end{itemize}

また、以下の点に注意してください:
\begin{itemize}
\item 確実に実在する論文のみを引用してください
\item 記憶が曖昧な場合は「要確認」と明記してください
\item 出典を特定できない主張は、「一般的な見解」と注記してください
\end{itemize}
\end{promptbox}

\begin{promptbox}
\textbf{【「わからない」を許可するプロンプト例】}

以下の質問に回答してください。ただし、以下のルールに従ってください:
\begin{itemize}
\item 確信がある情報のみを回答してください
\item 不確実な情報には [要確認] タグを付けてください
\item わからないことは「この点についてはわかりません」と\\
正直に回答してください
\item 推測で回答する場合は、「推測ですが」と明記してください
\end{itemize}

決して架空の情報を生成しないでください。
\end{promptbox}

\begin{screen}
\textbf{キーポイント}: ハルシネーション対策の本質は、「AIの出力を無条件に信頼しない」という姿勢である。これは「AIを使わない」ということではなく、「AIの出力を仮説として扱い、検証してから採用する」ということである。この姿勢は、科学的な研究態度そのものとも言える。
\end{screen}

%=========================================
\section{研究活動におけるプロンプト設計の実践}
\label{sec:2.6}
%=========================================

ここまで学んだ技法を、研究活動の具体的な場面に適用する方法を示す。

\subsection{文献調査用プロンプトの設計}

文献調査は、研究の初期段階で最もAIの支援が有効な場面である。

\begin{promptbox}
\textbf{【文献調査用プロンプト：サーベイの構造化】}

\textbf{[Role]} あなたは情報工学分野の研究に精通した\\
図書館司書です。学術文献の体系的な整理と分類を専門と\\
しています。

\textbf{[Objective]} 以下の研究テーマに関連する研究の\\
分類体系を提案してください。

\textbf{[Context]} 研究テーマ: 「大規模言語モデルの\\
ファインチューニング手法」 / 目的: 修士論文の関連研究\\
セクションの構成を決めるため / 対象期間: 2020年〜2025年

\textbf{[Format]} 大分類 $>$ 中分類 $>$ 小分類。\\
各分類に対して: 定義（50字程度）、代表的な手法名を\\
2〜3個、この分類が重要な理由

注意: 実在する手法名のみを使用してください。
\end{promptbox}

\textbf{ベストプラクティス:}
\begin{itemize}
\item 最初は広く探索し、徐々に焦点を絞る
\item AIが提示した論文は必ずGoogle ScholarやSemantic Scholarで確認する
\item サーベイの構造はAIに提案させ、個々の論文の評価は自分で行う
\end{itemize}

\textbf{アンチパターン:}
\begin{itemize}
\item AIが生成した文献リストをそのまま使う（ハルシネーションのリスク）
\item AIの要約だけで論文を読んだ気になる（原典の精読は不可欠）
\item 1回のプロンプトで完璧なサーベイを期待する（段階的に進める）
\end{itemize}

\subsection{論文執筆用プロンプトの設計}

\begin{promptbox}
\textbf{【論文執筆用プロンプト：Abstractの推敲】}

\textbf{[Role]} あなたは英語ネイティブの学術論文校正者です。\\
計算機科学分野の論文を年間100本以上校正しています。

\textbf{[Objective]} 以下の英語Abstractを推敲してください。\\
3つのバージョンを提示し、それぞれの特徴を説明してください。

\textbf{[Context]} 投稿先: ACL 2027（自然言語処理のトップ\\
会議） / 文字数制限: 250語以内 / 論文の主要な貢献:\\
新しいファインチューニング手法の提案とその有効性の実証

\textbf{[Format]} 各バージョンについて:\\
(1) 修正後のAbstract全文、\\
(2) 主な変更点の一覧（変更前→変更後）、\\
(3) そのバージョンの長所と短所
\end{promptbox}

\textbf{ベストプラクティス:}
\begin{itemize}
\item 最初に構成（アウトライン）を固め、その後で各セクションを執筆する
\item AIに生成させた文章は必ず自分で読み直し、加筆修正する
\item 専門用語の使い方が正しいか確認する
\end{itemize}

\textbf{アンチパターン:}
\begin{itemize}
\item 論文全体をAIに一括で書かせる（品質が低下し、一貫性も失われる）
\item AIの文体をそのまま使う（自分の文体で書き直すべき）
\item 投稿先のスタイルガイドを無視する
\end{itemize}

\subsection{コード生成用プロンプトの設計}

\begin{promptbox}
\textbf{【コード生成用プロンプト：データ分析スクリプト】}

\textbf{[Role]} あなたはPythonのデータ分析に精通したシニアエンジニアです。

\textbf{[Objective]} 以下の仕様に基づいて、データ分析用のPythonスクリプトを作成してください。

\textbf{[Context]} データ: CSVファイル（列: id, age, gender,\\
score1, score2, score3, group） / 行数: 約500行 /\\
目的: グループ間のスコア比較 /\\
環境: Python 3.11, pandas, scipy, matplotlib

\textbf{[仕様]}
\begin{enumerate}
\item CSVファイルを読み込む
\item 基本統計量（平均、標準偏差、中央値）をグループ別に算出
\item グループ間の差の検定（t検定またはMann-Whitney U検定をデータの正規性に基づいて自動選択）
\item 結果を可視化（箱ひげ図、棒グラフ）
\item 結果を\LaTeX 形式の表として出力
\end{enumerate}
\end{promptbox}

\textbf{ベストプラクティス:}
\begin{itemize}
\item 仕様を明確に記述し、あいまいさを排除する
\item 生成されたコードは必ず自分で動作確認する
\item テストケースもAIに生成させ、網羅的にテストする
\item コードの各行が何をしているか理解してから使用する
\end{itemize}

\textbf{アンチパターン:}
\begin{itemize}
\item 理解していないコードをそのまま使う（バグの温床）
\item エラーが出た際に、エラーメッセージを読まずにAIに丸投げする
\item セキュリティに関わるコード（認証、暗号化など）をAIに任せる
\end{itemize}

\subsection{プロンプトテンプレートの管理}

研究活動で頻繁に使うプロンプトは、テンプレートとして保存・管理することを推奨する。

\begin{lstlisting}[language={}, caption=推奨フォルダ構成]
prompts/
  literature_review/
    survey_structure.md     # サーベイの構造化
    paper_summary.md        # 論文の要約
    comparison_table.md     # 論文の比較表作成
  writing/
    abstract_review.md      # Abstract推敲
    related_work.md         # 関連研究セクション
    english_editing.md      # 英語校正
  coding/
    data_analysis.md        # データ分析
    debug.md                # デバッグ支援
    test_generation.md      # テスト生成
  presentation/
    slide_outline.md        # スライド構成
    qa_preparation.md       # 質疑応答準備
\end{lstlisting}

テンプレートの各ファイルには、以下の情報を含めるとよい:
\begin{enumerate}
\item テンプレートの目的と適用場面
\item プロンプトの全文（変数部分は \texttt{[ここに入力]} で示す）
\item 使用時の注意点
\item 過去の使用実績と改善履歴
\end{enumerate}

\begin{screen}
\textbf{ヒント}: プロンプトテンプレートをGitで管理すると、改善の履歴を追跡できる。チーム（研究室）で共有すれば、メンバー全員のAI活用スキルの底上げにもつながる。
\end{screen}

%=========================================
\section*{演習問題}
\addcontentsline{toc}{section}{演習問題}
%=========================================

\begin{enumerate}
\item \textbf{問題2.1（ROCFモデルの適用）}

以下のタスクに対して、ROCFモデルの4要素（Role, Objective, Context, Format）を明示的に含むプロンプトを設計せよ。

タスク: あなたの研究テーマ（または興味のあるテーマ）に関連する先行研究を5本リストアップし、各論文の手法と結果を比較する表を作成したい。

\item \textbf{問題2.2（プロンプト技法の使い分け）}

以下の3つのタスクに対して、それぞれ最適なプロンプト技法（Zero-shot, Few-shot, Chain-of-Thought）を選択し、その技法を用いたプロンプトを作成せよ。なぜその技法を選んだかも説明せよ。

\begin{enumerate}
\item[(a)] 英語の論文タイトルを日本語に翻訳する
\item[(b)] 研究論文のAbstractを「新規性」「有用性」「信頼性」の3観点で評価する
\item[(c)] 2つの研究手法の違いを分析し、どちらが特定の問題設定に適しているか判断する
\end{enumerate}

\item \textbf{問題2.3（ハルシネーション検出の実践）}

任意の対話型AIを使い、以下を実行せよ。

\begin{enumerate}
\item 自分の研究分野に関する質問をし、AIに論文を5本推薦させる
\item 推薦された5本の論文すべてについて、Google ScholarまたはSemantic Scholarで実在を確認する
\item 実在しない論文があった場合、その割合を記録する
\item ハルシネーションを減らすためにプロンプトを改善し、再度実行する
\item 改善前後の結果を比較し、レポートにまとめる
\end{enumerate}

\item \textbf{問題2.4（段階的精緻化の実践）}

任意の対話型AIを使い、以下の手順で段階的精緻化を実践せよ。

\begin{enumerate}
\item 自分の研究テーマのAbstract（200〜300字、日本語）のドラフトをAIに生成させる
\item 生成されたドラフトの問題点を3つ以上指摘し、AIに修正を依頼する
\item 修正後のバージョンをさらに改善するため、別の観点から修正を依頼する
\item 最終バージョンと最初のドラフトを比較し、どのように改善されたかを分析する
\end{enumerate}

\item \textbf{問題2.5（メタプロンプティング）}

AIに「効果的なプロンプトを設計してもらう」というメタプロンプティングを実践せよ。

\begin{enumerate}
\item 以下のタスクをAIに伝える: 「修士論文の中間発表用スライド（10枚程度）の構成案を作成したい」
\item AIに、このタスクに最適なプロンプトを設計させる
\item AIが設計したプロンプトを実際に使用し、結果を評価する
\item AIが設計したプロンプトの良い点と改善すべき点を分析する
\end{enumerate}

\item \textbf{問題2.6（総合演習：プロンプトの比較実験）}

同じタスクに対して、3種類のプロンプト（シンプルなプロンプト、ROCFモデルを適用したプロンプト、Few-shot + Chain-of-Thoughtを組み合わせたプロンプト）を作成し、それぞれの出力を比較せよ。

タスク: 「あなたの研究分野における、今後5年間の主要な研究トレンドを予測する」

評価の観点: 回答の具体性、回答の正確性（検証可能な情報の割合）、回答の有用性（研究活動に実際に役立つか）、ハルシネーションの有無
\end{enumerate}
